% GENERATED FILE. Edit canonical lessons, then run npm run generate:books.
% canonical-source-hash: fnv1a64:d4968f805cf4ad05
% canonical-lessons: IT-C331-ingredienti, IT-C331-dose, IT-C331-pizzico, IT-C331-cucchiaino, IT-C331-ciotola

\chapter{Gli ingredienti}
\label{ch:it-gli-ingredienti}
\hlchaptermodality{331}

\begin{chapteropening}
\textbf{By the end of this chapter:} \emph{I can read out a recipe: the ingredients, the amounts, a pinch, a teaspoon and a bowl.}

Complete the last lesson of chapter 331: I can read out a recipe: the ingredients, the amounts, a pinch, a teaspoon and a bowl.
\end{chapteropening}

\section[gli ingredienti]{gli ingredienti — the ingredients}

\label{lesson:IT-C331-ingredienti}



\begin{quote}
Before the new one: say the Italian for the excavations, then the Italian for the archaeologist.
\end{quote}

\subsection*{You'll want to know: gli ingredienti}
\textbf{gli ingredienti} — \textquotedblleft{}the ingredients\textquotedblright{}.

\begin{quote}
Gli ingredienti sono pochi: farina, uova e latte.
\end{quote}

\begin{quote}
\emph{There are only a few ingredients: flour, eggs and milk.}
\end{quote}

Singular \emph{l'ingrediente}.

\begin{cousinweb}
\emph{ingrediente} is a learned borrowing from Latin \emph{ingrediēns}, the present participle of \emph{ingredī} \textquotedblleft{}to go in\textquotedblright{} (Lewis \& Short): what goes into the dish.
\end{cousinweb}

\begin{grammarlens}[title={What you've built}]
One more word for telling, explaining and describing at B1.
\end{grammarlens}

\subsection*{Your turn}
\begin{itemize}
\raggedright
  \item \emph{Say it:} \emph{gli ingredienti}
  \item \emph{Say it:} \emph{gli ingredienti}, once more
  \item \emph{Say it:} a sentence of your own with \emph{gli ingredienti}
  \item \emph{Recall:} say the Italian for the excavations, then the Italian for the archaeologist, then say \emph{gli ingredienti} again
\end{itemize}

\begin{tcolorbox}[breakable,colback=teal!4,colframe=teal!35!black,title={Before you move on}]
What does gli ingredienti mean? (\textquotedblleft{}The ingredients\textquotedblright{}.) Say it once more.
\end{tcolorbox}
\section[la dose]{la dose — the amount, the quantity (in a recipe); the dose}

\label{lesson:IT-C331-dose}



\begin{quote}
Before the new one: say the Italian for the archaeologist, then the Italian for the ingredients.
\end{quote}

\subsection*{You'll want to know: la dose}
\textbf{la dose} — \textquotedblleft{}the amount, the quantity (in a recipe); the dose\textquotedblright{}.

\begin{quote}
Le dosi sono per quattro persone.
\end{quote}

\begin{quote}
\emph{The amounts are for four people.}
\end{quote}

In recipes: \emph{dosi per quattro}.

\begin{cousinweb}
\emph{dose} is usually traced to Greek \emph{dósis} \textquotedblleft{}a giving\textquotedblright{}, through Medieval Latin \emph{dosis} (EtymDB has no Italian entry).
\end{cousinweb}

\begin{grammarlens}[title={What you've built}]
One more word for telling, explaining and describing at B1.
\end{grammarlens}

\subsection*{Your turn}
\begin{itemize}
\raggedright
  \item \emph{Say it:} \emph{la dose}
  \item \emph{Say it:} \emph{la dose}, once more
  \item \emph{Say it:} a sentence of your own with \emph{la dose}
  \item \emph{Recall:} say the Italian for the archaeologist, then the Italian for the ingredients, then say \emph{la dose} again
\end{itemize}

\begin{tcolorbox}[breakable,colback=teal!4,colframe=teal!35!black,title={Before you move on}]
What does la dose mean? (\textquotedblleft{}The amount, the quantity (in a recipe); the dose\textquotedblright{}.) Say it once more.
\end{tcolorbox}
\section[il pizzico]{il pizzico — the pinch}

\label{lesson:IT-C331-pizzico}



\begin{quote}
Before the new one: say the Italian for the ingredients, then the Italian for the amount.
\end{quote}

\subsection*{You'll want to know: il pizzico}
\textbf{il pizzico} — \textquotedblleft{}the pinch\textquotedblright{}.

\begin{quote}
Aggiungi un pizzico di sale.
\end{quote}

\begin{quote}
\emph{Add a pinch of salt.}
\end{quote}

Built on \emph{pizzicare}, ‘to pinch’.

\begin{cousinweb}
\emph{pizzico} is built on \emph{pizzicare} \textquotedblleft{}to pinch\textquotedblright{} (EtymDB), a verb usually traced to the sound and gesture of pinching.
\end{cousinweb}

\begin{grammarlens}[title={What you've built}]
One more word for telling, explaining and describing at B1.
\end{grammarlens}

\subsection*{Your turn}
\begin{itemize}
\raggedright
  \item \emph{Say it:} \emph{il pizzico}
  \item \emph{Say it:} \emph{il pizzico}, once more
  \item \emph{Say it:} a sentence of your own with \emph{il pizzico}
  \item \emph{Recall:} say the Italian for the ingredients, then the Italian for the amount, then say \emph{il pizzico} again
\end{itemize}

\begin{tcolorbox}[breakable,colback=teal!4,colframe=teal!35!black,title={Before you move on}]
What does il pizzico mean? (\textquotedblleft{}The pinch\textquotedblright{}.) Say it once more.
\end{tcolorbox}
\section[il cucchiaino]{il cucchiaino — the teaspoon}

\label{lesson:IT-C331-cucchiaino}



\begin{quote}
Before the new one: say the Italian for the amount, then the Italian for the pinch.
\end{quote}

\subsection*{You'll want to know: il cucchiaino}
\textbf{il cucchiaino} — \textquotedblleft{}the teaspoon\textquotedblright{}.

\begin{quote}
Un cucchiaino di zucchero, per favore.
\end{quote}

\begin{quote}
\emph{A teaspoon of sugar, please.}
\end{quote}

A little \emph{cucchiaio}, ‘spoon’.

\begin{cousinweb}
\emph{cucchiaino} is a little \emph{cucchiaio} \textquotedblleft{}spoon\textquotedblright{}, from Latin \emph{cochlear} \textquotedblleft{}a spoon\textquotedblright{}, first a tool for eating snails, \emph{cochlea} (Lewis \& Short).
\end{cousinweb}

\begin{grammarlens}[title={What you've built}]
One more word for telling, explaining and describing at B1.
\end{grammarlens}

\subsection*{Your turn}
\begin{itemize}
\raggedright
  \item \emph{Say it:} \emph{il cucchiaino}
  \item \emph{Say it:} \emph{il cucchiaino}, once more
  \item \emph{Say it:} a sentence of your own with \emph{il cucchiaino}
  \item \emph{Recall:} say the Italian for the amount, then the Italian for the pinch, then say \emph{il cucchiaino} again
\end{itemize}

\begin{tcolorbox}[breakable,colback=teal!4,colframe=teal!35!black,title={Before you move on}]
What does il cucchiaino mean? (\textquotedblleft{}The teaspoon\textquotedblright{}.) Say it once more.
\end{tcolorbox}
\section[la ciotola]{la ciotola — the bowl}

\label{lesson:IT-C331-ciotola}



\begin{quote}
Before the new one: say the Italian for the pinch, then the Italian for the teaspoon.
\end{quote}

\subsection*{You'll want to know: la ciotola}
\textbf{la ciotola} — \textquotedblleft{}the bowl\textquotedblright{}.

\begin{quote}
Metti la farina in una ciotola.
\end{quote}

\begin{quote}
\emph{Put the flour in a bowl.}
\end{quote}

Also for pets: \emph{la ciotola del gatto}.

\begin{cousinweb}
\emph{ciotola} is usually traced to a little form of Latin \emph{cyathus} \textquotedblleft{}a ladle, a cup\textquotedblright{} (Lewis \& Short); the path is not settled.
\end{cousinweb}

\begin{grammarlens}[title={What you've built}]
One more word for telling, explaining and describing at B1.
\end{grammarlens}

\subsection*{Your turn}
\begin{itemize}
\raggedright
  \item \emph{Say it:} \emph{la ciotola}
  \item \emph{Say it:} \emph{la ciotola}, once more
  \item \emph{Say it:} a sentence of your own with \emph{la ciotola}
  \item \emph{Recall:} say the Italian for the pinch, then the Italian for the teaspoon, then say \emph{la ciotola} again
\end{itemize}

\begin{tcolorbox}[breakable,colback=teal!4,colframe=teal!35!black,title={Before you move on}]
What does la ciotola mean? (\textquotedblleft{}The bowl\textquotedblright{}.) Say it once more.
\end{tcolorbox}
